\subsection{Менеджер памяти-1}

\subsubsection{Описание решения}

В задаче необходимо реализовать прототип аллокатора. Воспользуемся красно-чёрными деревьями для хранения свободных блоков памяти в отсортированном виде. Каждый блок памяти представим в виде пары значений: начального адреса в памяти и его размера.

Пусть блоки памяти будут упорядочены во множестве по размеру, а в ассоциативном массиве --- по начальному адресу. В момент запуска программы вся выделенная память доступна для аллокаций.

Если поступил запрос на аллокацию блока размера \( K \), ищем во множестве минимальный по размеру свободный блок, который вмещает в себя \( K \) ячеек. Если такой не найден, то выводим ошибку аллокации. Иначе обновляем размер и начальный адрес выбранного свободного блока и вставляем обратно \textit{(если осталось свободное место)}. Аллоцированный блок данных сохраняем в хеш-таблицу по связке с номером запроса.

Если поступил запрос на деаллокацию блока, аллоцированного в \( i \)-м запросе, достаём его из хеш-таблицы и удаляем запись. Возвращаем его в ассоциативный массив и смотрим на соседние блоки. Если они примыкают вплотную, то производим \textit{дефрагментацию}, то есть сливаем их в единый блок. Попутно синхронизируем записи во множестве свободных блоков.

Симуляция аллокатора завершена.

\subsubsection{Асимптотическая сложность}

\textit{Временная сложность решения} --- линейно-логарифмическая \( \mathcal{O}(M\cdot\log M) \). В случае аллокации поиск наименьшего свободного блока памяти во множестве занимает логарифмическое время. Удаление и вставка вершин в красно-чёрные деревья тоже имеют логарифмическую сложность.

В случае деаллокации поиск выделенного блока по номеру запроса в среднем выполняется за константное время. Другие используемые операции над красно-чёрными деревьями выполняются последовательно и за логарифмическое время.

\textit{Пространственная сложность решения} --- линейная \( \mathcal{O}(M) \). В памяти хранятся данные о выделенных блоках в нескольких структурах для оптимизации времени доступа. Этих записей не может быть больше запросов на аллокацию блоков в каждой из структур. Промежуточные данные хранятся в константном количестве памяти.

\subsubsection{Листинг кода}

\begin{lstlisting}
#include <cstddef>
#include <cstdint>
#include <cstdlib>
#include <iostream>
#include <iterator>
#include <map>
#include <set>
#include <unordered_map>

namespace {
struct Block {
  std::size_t start;
  std::size_t size;

  bool operator<(const Block& other) const {
    if (size != other.size) {
      return size < other.size;
    }
    return start < other.start;
  }
};

void Allocate(
    std::set<Block>& free_heap,
    std::map<std::size_t, std::size_t>& free_blocks,
    std::unordered_map<std::size_t, Block>& query_map,
    const std::size_t alloc_size,
    const std::size_t query_id
) {
  auto iter = free_heap.lower_bound({0, alloc_size});
  if (iter != free_heap.end()) {
    const Block block = *iter;
    free_heap.erase(iter);
    free_blocks.erase(block.start);

    const std::size_t remaining_size = block.size - alloc_size;
    if (remaining_size > 0) {
      const std::size_t new_start = block.start + alloc_size;
      free_heap.insert({new_start, remaining_size});
      free_blocks[new_start] = remaining_size;
    }
    query_map[query_id] = {block.start, alloc_size};
    std::cout << block.start + 1 << "\n";
  } else {
    std::cout << 1 << "\n";
  }
}

void Deallocate(
    std::set<Block>& free_heap,
    std::map<std::size_t, std::size_t>& free_blocks,
    std::unordered_map<std::size_t, Block>& query_map,
    const std::size_t query_id
) {
  auto dict_it = query_map.find(query_id);
  if (dict_it != query_map.end()) {
    const Block allocated_block = dict_it->second;
    query_map.erase(dict_it);

    auto insert_res = free_blocks.insert({allocated_block.start, allocated_block.size});
    auto curr_it = insert_res.first;

    auto next_it = std::next(curr_it);
    if (next_it != free_blocks.end() &&
        allocated_block.start + allocated_block.size == next_it->first) {
      free_heap.erase({next_it->first, next_it->second});

      curr_it->second += next_it->second;
      free_blocks.erase(next_it);
    }

    if (curr_it != free_blocks.begin()) {
      auto prev_it = std::prev(curr_it);
      if (prev_it->first + prev_it->second == curr_it->first) {
        free_heap.erase({prev_it->first, prev_it->second});

        prev_it->second += curr_it->second;
        free_blocks.erase(curr_it);

        curr_it = prev_it;
      }
    }
    free_heap.insert({curr_it->first, curr_it->second});
  }
}
}  // namespace

int main() {
  std::size_t n_count = 0;
  std::size_t m_count = 0;
  std::cin >> n_count >> m_count;

  std::set<Block> free_heap;
  std::map<std::size_t, std::size_t> free_blocks;
  std::unordered_map<std::size_t, Block> query_map;
  free_heap.insert({0, n_count});
  free_blocks[0] = n_count;

  for (std::size_t i = 0; i < m_count; i++) {
    int32_t item = 0;
    std::cin >> item;
    if (item > 0) {
      const auto alloc_size = static_cast<std::size_t>(item);
      Allocate(free_heap, free_blocks, query_map, alloc_size, i + 1);
      continue;
    }
    const std::size_t query_id = std::abs(item);
    Deallocate(free_heap, free_blocks, query_map, query_id);
  }
  return 0;
}
\end{lstlisting}